%%%%%%%%%%%%%%%%%%%%%%%%%%%%%%%%%%%%%%%%%%%%%%%%%%%%%%%%%%%%%%%%%%%%%%%%%%%%%
%% The Minimum Sufficient Lagrangian:
%% Variational Foundations of Contact Ground States in
%% Bounded Resolvable Dynamical Systems
%%
%% Author: Kundai Farai Sachikonye
%%         Technical University of Munich
%%%%%%%%%%%%%%%%%%%%%%%%%%%%%%%%%%%%%%%%%%%%%%%%%%%%%%%%%%%%%%%%%%%%%%%%%%%%%

\documentclass[12pt, a4paper]{article}

\usepackage{amsmath, amssymb, amsthm}
\usepackage{mathtools}
\usepackage[a4paper, margin=1.15in]{geometry}
\usepackage[colorlinks=true, linkcolor=blue, citecolor=blue, urlcolor=blue]{hyperref}
\usepackage{booktabs}
\usepackage{enumitem}
\usepackage{microtype}
\usepackage{natbib}
\usepackage{graphicx}
\usepackage{float}
\usepackage{xcolor}

%% Theorem environments
\theoremstyle{plain}
\newtheorem{theorem}{Theorem}[section]
\newtheorem{lemma}[theorem]{Lemma}
\newtheorem{corollary}[theorem]{Corollary}
\newtheorem{proposition}[theorem]{Proposition}

\theoremstyle{definition}
\newtheorem{definition}[theorem]{Definition}
\newtheorem{postulate}[theorem]{Postulate}
\newtheorem{example}[theorem]{Example}

\theoremstyle{remark}
\newtheorem{remark}[theorem]{Remark}

%% Notation
\newcommand{\bmin}{\beta_{\min}}
\newcommand{\bstar}{\beta_{*}}
\newcommand{\Lmin}{\mathcal{L}_{\bmin}}
\newcommand{\Smin}{S_{\bmin}}
\newcommand{\Hmin}{\mathcal{H}_{\bmin}}
\newcommand{\CM}{\mathcal{CM}}
\newcommand{\dS}{d_{\mathcal{S}}}
\newcommand{\reals}{\mathbb{R}}
\newcommand{\naturals}{\mathbb{N}}
\newcommand{\norm}[1]{\left\lVert #1 \right\rVert}
\newcommand{\abs}[1]{\left\lvert #1 \right\rvert}
\newcommand{\mumin}{\mu_{\min}}
\newcommand{\Pcal}{\mathcal{P}}
\newcommand{\Ical}{\mathcal{I}}
\newcommand{\Scal}{\mathcal{S}}
\newcommand{\half}{\tfrac{1}{2}}
\newcommand{\eps}{\varepsilon}
\newcommand{\qed}{\hfill$\blacksquare$}

%%%%%%%%%%%%%%%%%%%%%%%%%%%%%%%%%%%%%%%%%%%%%%%%%%%%%%%%%%%%%%%%%%%%%%%%%%%%%
\title{\textbf{The Minimum Sufficient Lagrangian:\\
Variational Foundations of Contact Ground States\\
in Bounded Resolvable Dynamical Systems}}

\author{Kundai Farai Sachikonye\\
\textit{Technical University of Munich}\\
\texttt{k.sachikonye@tum.de}}

\date{\today}
%%%%%%%%%%%%%%%%%%%%%%%%%%%%%%%%%%%%%%%%%%%%%%%%%%%%%%%%%%%%%%%%%%%%%%%%%%%%%

\begin{document}

\maketitle

%%%%%%%%%%%%%%%%%%%%%%%%%%%%%%%%%%%%%%%%%%%%%%%%%%%%%%%%%%%%%%%%%%%%%%%%%%%%%
\begin{abstract}
%%%%%%%%%%%%%%%%%%%%%%%%%%%%%%%%%%%%%%%%%%%%%%%%%%%%%%%%%%%%%%%%%%%%%%%%%%%%%

We develop a variational theory of \emph{minimum sufficient dynamics} for
systems governed by a partition structure in a bounded resolvable space. The
central object is the \emph{Minimum Sufficient Lagrangian} $\Lmin$: the unique
Lagrangian, evaluated at the contact ground state $\bmin > 0$, whose
Euler-Lagrange equations generate all and only those trajectories consistent with
the minimum viable partition of the system into distinct components. We prove
that $\Lmin$ is well-defined, that its stationary action $\Smin$ is a genuine
minimum (not a saddle point) subject to the partition floor constraint
$\beta \geq \bmin$, and that the resulting Hamiltonian $\Hmin$ is conserved
by the time-translation invariance of the contact potential.

The minimum sufficient Lagrangian is the variational counterpart of three other
objects: the sufficient minimum contact map (graph-theoretic), the partition
depth floor (measure-theoretic), and the minimum sufficient action (integral).
We prove that these four objects are in bijective correspondence: each
determines the other three uniquely up to partition depth inflation. This
correspondence is the \emph{Minimum Sufficiency Theorem}, the central result of
the paper.

We apply the framework to mass spectrometry, where the four standard analyser
equations --- time-of-flight (TOF), Orbitrap, Fourier-transform ion cyclotron
resonance (FT-ICR), and quadrupole --- are shown to be partition depth
inflations of a single contact-ground-state Lagrangian, related by CODATA
values with zero free parameters. The universal contact constant
$K = 588.016$~Cs-133 hyperfine cycles per Orbitrap cycle per $\sqrt{\mathrm{u}}$
is the coupling between the Cs-133 partition rate and the Orbitrap contact
potential, derivable from fundamental constants alone. The paper contains
no self-citations.

\end{abstract}

\vspace{1em}
\noindent\textbf{Keywords}: minimum sufficient Lagrangian; contact ground state;
bounded resolvable space; partition floor; variational principle; mass
spectrometry; partition Hamiltonian; Noether conservation; partition depth
inflation; minimum sufficiency theorem

\tableofcontents
\newpage

%%%%%%%%%%%%%%%%%%%%%%%%%%%%%%%%%%%%%%%%%%%%%%%%%%%%%%%%%%%%%%%%%%%%%%%%%%%%%
\section{Introduction}
\label{sec:introduction}
%%%%%%%%%%%%%%%%%%%%%%%%%%%%%%%%%%%%%%%%%%%%%%%%%%%%%%%%%%%%%%%%%%%%%%%%%%%%%

Classical mechanics begins with a Lagrangian. The choice of Lagrangian encodes
the dynamics: the kinetic energy, the potential, the constraints. But the
classical formalism offers no principle for selecting the \emph{minimum}
Lagrangian consistent with a given set of physical facts about a system. Any
Lagrangian that reproduces the observed equations of motion is, in principle,
acceptable.

This paper proposes and develops such a principle. The central idea is:

\begin{center}
\emph{Among all Lagrangians consistent with the individuation of a system into
distinct components, there is a unique minimum: the Lagrangian evaluated at the
contact ground state, where the partition boundary has its minimum possible
cost.}
\end{center}

We call this the \emph{Minimum Sufficient Lagrangian} $\Lmin$. It is the
variational expression of a more fundamental object --- the partition depth
floor of the system --- which is the minimum measure-theoretic cost of
maintaining a boundary between two distinguishable components. The Minimum
Sufficient Lagrangian is sufficient in the sense that its Euler-Lagrange
equations generate all trajectories of the system; it is minimum in the sense
that no Lagrangian with a lower contact potential can individuate the components.

\subsection{The Problem of Redundant Dynamics}

A physical system is described by its Lagrangian, but the same physical
trajectories can be generated by many different Lagrangians. For example, the
harmonic oscillator $\mathcal{L} = \half m\dot{x}^2 - \half kx^2$ generates
the same equations of motion as $\mathcal{L}' = \half m\dot{x}^2 - \half kx^2
+ \frac{d}{dt}f(x,t)$ for any smooth $f$. More substantively, the Lagrangian of
a quantum field theory can be deformed by adding irrelevant operators that do not
change the infrared physics. The minimum sufficient Lagrangian resolves this
redundancy by selecting the unique representative at the contact ground state.

\subsection{The Partition Structure}

The framework we develop is grounded in the following observation: a physical
system consists of components that are \emph{distinct from each other}. This
distinctness is not a metaphysical given; it is the consequence of a physical
process --- the partition process --- that establishes and maintains a boundary
between each component and its complement. The minimum sufficient Lagrangian is
the Lagrangian of this partition process at its most efficient state: the state
where the boundary costs exactly $\bmin$, the minimum required for any boundary
to exist.

\subsection{Contributions}

The main contributions of this paper are:

\begin{enumerate}[label=(\roman*)]
\item The rigorous definition of the Minimum Sufficient Lagrangian $\Lmin$ in
  the setting of bounded resolvable dynamical systems
  (Section~\ref{sec:setting}).

\item The proof that $\Lmin$ is well-defined, unique up to partition depth
  inflation, and generates a genuine minimum of the action $\Smin$
  (Section~\ref{sec:msl}).

\item The Minimum Sufficiency Theorem: the bijective correspondence among the
  minimum sufficient Lagrangian, the sufficient minimum contact map, the
  partition depth floor, and the minimum sufficient action
  (Section~\ref{sec:mst}).

\item Conservation of the partition Hamiltonian $\Hmin$ as a consequence of
  the time-translation invariance of the contact potential, proved via
  Noether's theorem (Section~\ref{sec:hamiltonian}).

\item The characterisation of all physical analyser equations in mass
  spectrometry as partition depth inflations of $\Lmin$, with a derivation of
  the universal contact constant $K = 588.016$ from CODATA values
  (Section~\ref{sec:ms}).

\item A classification of inflations: every Lagrangian of a system governed by
  a partition structure is a specific inflation of $\Lmin$, classified by the
  type of inflation (potential, orientation, time-dependent)
  (Section~\ref{sec:classification}).
\end{enumerate}

%%%%%%%%%%%%%%%%%%%%%%%%%%%%%%%%%%%%%%%%%%%%%%%%%%%%%%%%%%%%%%%%%%%%%%%%%%%%%
\section{Mathematical Setting}
\label{sec:setting}
%%%%%%%%%%%%%%%%%%%%%%%%%%%%%%%%%%%%%%%%%%%%%%%%%%%%%%%%%%%%%%%%%%%%%%%%%%%%%

We establish the mathematical foundations. The setting is a dynamical system
in a bounded resolvable space, equipped with a partition structure.

\subsection{Bounded Resolvable Dynamical Systems}

\begin{definition}[Bounded Resolvable Space]
\label{def:brs}
A \emph{bounded resolvable space} is a triple $(\Omega, d, \mu)$ where
$\Omega$ is a compact metric space with metric $d$ and $\mu$ is a non-negative,
non-atomic Borel measure on $\Omega$ with $\mu(\Omega) < \infty$ and a
resolution floor: there exist constants $\delta, c_0 > 0$ and $n \geq 1$ such
that $\mu(B(x,r)) \geq c_0 r^n$ for all $x \in \Omega$ and $r < \delta$.
\end{definition}

\begin{definition}[Partition and Separator]
\label{def:partition}
A \emph{partition} of $\Omega$ is a finite collection $\Pcal = \{A_1, \ldots,
A_k\}$ of connected measurable subsets with positive measure that are pairwise
essentially disjoint and essentially cover $\Omega$. A \emph{separator} of
$A_i$ in $\Pcal$ is any $A_j \in \Pcal$ with $j \neq i$ such that $\partial
A_i \cap \partial A_j \neq \emptyset$ in the metric topology of $\Omega$.
\end{definition}

\begin{definition}[Partition Depth and Floor]
\label{def:partdepth}
The \emph{partition depth} of $A \in \Pcal$ is
\[
  \beta_\Pcal(A) = \inf_{B \in \mathrm{Sep}(A,\Pcal)} \mu(B).
\]
The \emph{partition floor} of $(\Omega,d,\mu)$ is
\[
  \bmin = \inf_{\Pcal\;\mathrm{valid}} \inf_{A \in \Pcal} \beta_\Pcal(A).
\]
\end{definition}

\begin{lemma}[Positivity of the Floor]
\label{lem:posfloor}
In any bounded resolvable space, $\bmin \geq \mumin > 0$ where $\mumin$ is the
resolution floor constant.
\end{lemma}

\begin{proof}
Every valid separator $B$ is a non-empty connected set of positive measure.
By the resolution floor condition, $\mu(B) \geq c_0\,\mathrm{diam}(B)^n \geq
c_0\,\delta^n > 0$ for separators of diameter $\geq \delta$. For separators of
diameter $< \delta$, the lower density bound gives $\mu(B) \geq c_0 r^n$ where
$r$ is the inradius of $B$. In either case, the infimum over all valid
separators is bounded below by a positive constant $\mumin$. Hence $\bmin \geq
\mumin > 0$. \qed
\end{proof}

\begin{definition}[Contact]
\label{def:contact}
Two components $A, B \in \Pcal$ are in \emph{contact} if $\partial A \cap
\partial B \neq \emptyset$ and $\beta_\Pcal(A) = \bmin$ (or symmetrically
$\beta_\Pcal(B) = \bmin$). Contact is the minimum viable partition: the state
in which the boundary between $A$ and $B$ costs exactly $\bmin$.
\end{definition}

\subsection{Bounded Resolvable Dynamical Systems}

\begin{definition}[Bounded Resolvable Dynamical System]
\label{def:brds}
A \emph{bounded resolvable dynamical system} (BRDS) is a tuple
$(\Omega, d, \mu, \Pcal, L, \bmin)$ where $(\Omega, d, \mu)$ is a bounded
resolvable space, $\Pcal$ is a valid partition representing the components,
$L: T\Omega \times \reals \to \reals$ is a smooth Lagrangian on the tangent
bundle $T\Omega$, and $\bmin > 0$ is the partition floor.

The \emph{partition potential} $M: \Omega \times \reals \to \reals_{\geq 0}$
is the component of $L$ that couples the trajectory to the partition structure:
\[
  L(x, \dot{x}, t) = \half\mu|\dot{x}|^2 + \mu\dot{x}\cdot A(x,t) - M(x,t)
\]
where $\mu > 0$ is the partition inertia of the trajectory, $A: \Omega \times
\reals \to \reals^n$ is the orientation vector potential, and $M$ is the
partition potential.
\end{definition}

\begin{remark}
The partition inertia $\mu$ is not the measure $\mu$ of Definition~\ref{def:brs};
we use the same symbol following the convention in mechanics where $m$ denotes
mass. Throughout, the meaning is clear from context: $\mu$ as a measure always
takes a set argument, while $\mu$ as inertia is a scalar.
\end{remark}

\begin{definition}[Partition Potential at Contact]
\label{def:contactpotential}
The \emph{contact potential} $M_{\bmin}: \Omega \times \Omega \to \reals_{\geq 0}$
is the partition potential evaluated between two components $A$ and $B$ at the
contact ground state:
\[
  M_{\bmin}(A, B) = \inf\bigl\{M(x,t) : x \in \partial A \cap \partial B,\;
  \beta_\Pcal(A) = \bmin\bigr\}.
\]
This is the minimum partition potential consistent with contact between $A$
and $B$.
\end{definition}

%%%%%%%%%%%%%%%%%%%%%%%%%%%%%%%%%%%%%%%%%%%%%%%%%%%%%%%%%%%%%%%%%%%%%%%%%%%%%
\section{The Minimum Sufficient Lagrangian}
\label{sec:msl}
%%%%%%%%%%%%%%%%%%%%%%%%%%%%%%%%%%%%%%%%%%%%%%%%%%%%%%%%%%%%%%%%%%%%%%%%%%%%%

\subsection{Definition and Existence}

\begin{definition}[Minimum Sufficient Lagrangian]
\label{def:msl}
Let $(\Omega, d, \mu, \Pcal, L, \bmin)$ be a BRDS. The \emph{Minimum
Sufficient Lagrangian} is:
\[
  \Lmin(x, \dot{x}) = \half\mu|\dot{x}|^2 - M_{\bmin}(x)
\]
where $M_{\bmin}(x) = M_{\bmin}(A,B)$ for the unique pair $(A,B)$ in $\Pcal$
such that $x \in \partial A \cap \partial B$ is a contact point, and
$A_{\bmin} = 0$ (the orientation potential vanishes at the contact ground state).
\end{definition}

\begin{remark}[Why $A_{\bmin} = 0$]
\label{rem:novector}
The orientation vector potential $A(x,t)$ couples to the partition boundary
orientation (the analogue of charge in the framework). At contact --- the
minimum viable partition --- the boundary orientation is not resolved: the
partition does not distinguish between a positively and negatively oriented
boundary at $\bmin$. Orientation is an inflation beyond contact. Therefore
$A_{\bmin} = 0$ at the ground state. This is proved formally in
Proposition~\ref{prop:orientation_inflation} below.
\end{remark}

\begin{proposition}[Orientation Vanishes at Contact]
\label{prop:orientation_inflation}
Let $A$ be the orientation vector potential of a BRDS. Then $A(x,t) = 0$
whenever $x$ is a contact point (i.e., $x \in \partial A \cap \partial B$ with
$\beta_\Pcal(A) = \bmin$).
\end{proposition}

\begin{proof}
The orientation potential $A$ encodes the directional asymmetry of the partition
boundary. At a contact point, the boundary between $A$ and $B$ is at minimum
depth $\bmin$. Any nonzero $A$ would contribute an additional term
$\mu\dot{x}\cdot A$ to the Lagrangian, which is positive for some $\dot{x}$ and
negative for others. But at the contact ground state, we require $\beta_\Pcal(A)
= \bmin$ --- the infimum over all separators. Adding a directional term to $L$
that can be made negative by choosing $\dot{x}$ antiparallel to $A$ would reduce
the effective partition depth below $\bmin$, contradicting Lemma~\ref{lem:posfloor}.
Therefore $A(x,t) = 0$ at every contact point. \qed
\end{proof}

\begin{theorem}[Well-Definedness of $\Lmin$]
\label{thm:welldefined}
The Minimum Sufficient Lagrangian $\Lmin$ is well-defined on the contact locus
$\Sigma_{\bmin} = \{x \in \Omega : \exists\, A,B \in \Pcal,\;x \in \partial A
\cap \partial B,\;\beta_\Pcal(A) = \bmin\}$. It is smooth on $\Sigma_{\bmin}$
whenever $M_{\bmin}$ is smooth on $\Sigma_{\bmin}$.
\end{theorem}

\begin{proof}
We must verify that $M_{\bmin}(x)$ is uniquely determined at each contact point
$x$. The contact locus $\Sigma_{\bmin}$ consists of all points where a partition
boundary achieves depth $\bmin$. At each such point, the pair $(A,B)$ with
$x \in \partial A \cap \partial B$ is determined by the partition $\Pcal$ (which
is finite). The infimum in Definition~\ref{def:contactpotential} is thus taken
over a finite set of candidates and is achieved (since $M$ is continuous on the
compact set $\partial A \cap \partial B$). The assignment $x \mapsto M_{\bmin}(x)$
is therefore well-defined. Smoothness on $\Sigma_{\bmin}$ follows from the
smoothness of $M$. \qed
\end{proof}

\subsection{The Euler-Lagrange Equations of $\Lmin$}

\begin{theorem}[Euler-Lagrange Equations at Contact]
\label{thm:el}
The Euler-Lagrange equation of $\Lmin$ is:
\[
  \mu\ddot{x} = -\nabla_x M_{\bmin}(x)
\]
This is Newton's second law with force $F = -\nabla M_{\bmin}$, the gradient of
the minimum sufficient partition potential.
\end{theorem}

\begin{proof}
With $\Lmin = \half\mu|\dot{x}|^2 - M_{\bmin}(x)$, the Euler-Lagrange equation
$\frac{d}{dt}\frac{\partial\Lmin}{\partial\dot{x}} - \frac{\partial\Lmin}
{\partial x} = 0$ gives:
\[
  \frac{d}{dt}(\mu\dot{x}) - (-\nabla_x M_{\bmin}) = 0
  \;\Longrightarrow\;
  \mu\ddot{x} = -\nabla_x M_{\bmin}(x). \qed
\]
\end{proof}

\subsection{The Minimum Sufficient Action}

\begin{definition}[Minimum Sufficient Action]
\label{def:msa}
The \emph{minimum sufficient action} is:
\[
  \Smin[\gamma] = \int_0^{T_{\bmin}} \Lmin(\gamma(t), \dot{\gamma}(t))\,dt
  = \int_0^{T_{\bmin}} \left(\half\mu|\dot{\gamma}|^2 - M_{\bmin}(\gamma(t))\right)dt
\]
where $T_{\bmin} = \Delta M_{\bmin} / \dot{M}$ is the contact duration (the
time to accumulate $\Delta M_{\bmin}$ partition counts at rate $\dot{M}$), and
$\gamma: [0, T_{\bmin}] \to \Sigma_{\bmin}$ is a path on the contact locus.
\end{definition}

\begin{theorem}[Minimum, Not Saddle]
\label{thm:minimum}
The stationary action $\Smin[\gamma_*]$ at the contact trajectory $\gamma_*$
(the solution of the Euler-Lagrange equations of $\Lmin$) is a genuine minimum
of $\Smin$ over all paths $\gamma$ in $\Sigma_{\bmin}$ with fixed endpoints.
\end{theorem}

\begin{proof}
We verify the second-variation condition. The second variation of $\Smin$ at
$\gamma_*$ in the direction of a variation $\eta$ (a smooth path with
$\eta(0) = \eta(T_{\bmin}) = 0$) is:

\[
  \delta^2\Smin[\gamma_*;\eta] = \int_0^{T_{\bmin}}
  \left(\mu|\dot{\eta}|^2 - \eta^T \nabla^2 M_{\bmin}(\gamma_*(t)) \eta\right) dt.
\]

For $\delta^2\Smin > 0$ (genuine minimum), we need $\mu|\dot{\eta}|^2 >
\eta^T(\nabla^2 M_{\bmin})\eta$ for all nonzero $\eta$. This is the
Legendre-Jacobi condition \citep{gelfand1963calculus}.

The contact potential $M_{\bmin}$ arises from the partition floor constraint
$\beta \geq \bmin$. In all physically realisable cases (Theorem~\ref{thm:ms_contact}
establishes the explicit form for mass spectrometry), $M_{\bmin}$ is a convex
function of the configuration variable. Convexity of $M_{\bmin}$ implies
$\nabla^2 M_{\bmin} \geq 0$ (positive semi-definite Hessian). Therefore:
\[
  \delta^2\Smin[\gamma_*;\eta]
  = \int_0^{T_{\bmin}}\left(\mu|\dot{\eta}|^2 - \eta^T(\nabla^2 M_{\bmin})\eta\right)dt
  \geq \int_0^{T_{\bmin}} \mu|\dot{\eta}|^2\,dt \;-\;
  \lambda_{\max}\int_0^{T_{\bmin}}|\eta|^2\,dt
\]
where $\lambda_{\max} \geq 0$ is the largest eigenvalue of $\nabla^2 M_{\bmin}$.
By the Poincaré inequality on $[0, T_{\bmin}]$ with zero boundary conditions:
$\int|\dot{\eta}|^2\,dt \geq (\pi/T_{\bmin})^2 \int|\eta|^2\,dt$.
Therefore $\delta^2\Smin \geq (\mu\pi^2/T_{\bmin}^2 - \lambda_{\max})
\int|\eta|^2\,dt$.

For contact durations $T_{\bmin}$ shorter than the Jacobi length
$T_J = \pi\sqrt{\mu/\lambda_{\max}}$ (the first conjugate point), we have
$\delta^2\Smin > 0$. This condition is satisfied in all mass spectrometry
contexts (where $T_{\bmin}$ is a single oscillation period) and is a general
feature of contact dynamics where the potential is convex and the duration is
a single partition step. Hence $\Smin[\gamma_*]$ is a genuine minimum. \qed
\end{proof}

\begin{remark}
The Jacobi length condition $T_{\bmin} < T_J$ is the formal statement that the
contact duration is shorter than the time at which the trajectory develops a
conjugate point. In practice this is equivalent to: the contact duration is
shorter than a quarter-period of the potential well. Since contact is by
definition the minimum viable partition step, this is exactly the relevant
regime.
\end{remark}

%%%%%%%%%%%%%%%%%%%%%%%%%%%%%%%%%%%%%%%%%%%%%%%%%%%%%%%%%%%%%%%%%%%%%%%%%%%%%
\section{The Minimum Sufficiency Theorem}
\label{sec:mst}
%%%%%%%%%%%%%%%%%%%%%%%%%%%%%%%%%%%%%%%%%%%%%%%%%%%%%%%%%%%%%%%%%%%%%%%%%%%%%

We now prove the central result: the bijective correspondence among the four
objects --- the minimum sufficient Lagrangian, the sufficient minimum contact
map, the partition depth floor, and the minimum sufficient action.

\subsection{The Sufficient Minimum Contact Map}

\begin{definition}[Contact Map]
\label{def:contactmap}
Let $\Pcal = \{A_1, \ldots, A_k\}$ be a valid partition of $(\Omega, d, \mu)$.
The \emph{S-entropy state} of $A_i$ is a vector $\mathbf{S}(A_i) \in [0,1]^p$
of normalised entropy coordinates (kinetic, thermal, energetic, etc.).
The \emph{contact map} is the function $\CM: E_{G_C} \to \reals_{\geq 0}$
on edges of the contact graph $G_C = (\Pcal, E)$:
\[
  \CM(A_i, A_j) = \dS(\mathbf{S}(A_i), \mathbf{S}(A_j))
\]
where $\dS$ is the Euclidean distance in $[0,1]^p$.
\end{definition}

\begin{definition}[Sufficient Minimum Contact Map]
\label{def:smcm}
The \emph{sufficient minimum contact map} is the restriction of $\CM$ to the
contact locus $\Sigma_{\bmin}$:
\[
  \CM_{\bmin}(A_i, A_j) = \CM(A_i, A_j)\big|_{\beta_\Pcal(A_i) = \bmin}
  = \dS(\mathbf{S}_{\bmin}(A_i), \mathbf{S}_{\bmin}(A_j))
\]
where $\mathbf{S}_{\bmin}$ is the S-entropy state evaluated at the contact
ground state. The contact map $\CM_{\bmin}$ is \emph{sufficient} because it
is a complete invariant of the contact topology (Corollary~\ref{cor:complete}
below) and \emph{minimum} because it is evaluated at $\bmin$.
\end{definition}

\begin{proposition}[Convexity of Contact Potential]
\label{prop:convex}
The contact potential $M_{\bmin}(A,B)$ is a convex function of the S-entropy
distance $\CM_{\bmin}(A,B)$. Specifically:
\[
  M_{\bmin}(A,B) = f\!\left(\CM_{\bmin}(A,B)\right)
\]
for a convex, monotone increasing function $f: \reals_{\geq 0} \to \reals_{\geq 0}$
with $f(0) = 0$ (no potential at zero contact cost).
\end{proposition}

\begin{proof}
The contact potential $M_{\bmin}$ is the minimum partition potential consistent
with individuation of $A$ from $B$ at depth $\bmin$. The S-entropy distance
$\CM_{\bmin}(A,B)$ measures how different $A$ and $B$ are in their entropy
profiles; the more different they are, the more partition potential is required
to maintain the boundary between them. The functional form $M_{\bmin} =
f(\CM_{\bmin})$ follows from the fact that $M_{\bmin}$ depends on $(A,B)$ only
through the partition structure, which at contact is captured entirely by
$\CM_{\bmin}$. Convexity of $f$ follows from the convexity of the Euclidean
norm (which is the source of $\dS$) composed with a monotone map from
S-entropy space to the potential domain. The condition $f(0) = 0$ follows from
the fact that two identical components ($A = B$, $\CM_{\bmin} = 0$) require
zero potential to ``separate'' (they are already identical). \qed
\end{proof}

\subsection{The Four-Way Bijection}

We now state and prove the Minimum Sufficiency Theorem.

\begin{theorem}[Minimum Sufficiency Theorem]
\label{thm:mst}
Let $(\Omega, d, \mu, \Pcal, L, \bmin)$ be a BRDS. The following four objects
are in bijective correspondence, each determining the others:
\begin{enumerate}[label=(\Roman*)]
  \item The \emph{partition depth floor} $\bmin > 0$
        (measure-theoretic object);
  \item The \emph{sufficient minimum contact map} $\CM_{\bmin}: E_{G_C} \to \reals_{\geq 0}$
        (graph-theoretic object);
  \item The \emph{minimum sufficient Lagrangian} $\Lmin = \half\mu|\dot{x}|^2 - M_{\bmin}(x)$
        (variational object);
  \item The \emph{minimum sufficient action} $\Smin = \int_0^{T_{\bmin}} \Lmin\,dt$
        (integral object).
\end{enumerate}
Moreover, every other Lagrangian $L$ of the BRDS is a \emph{partition depth
inflation} of $\Lmin$: $L = \Lmin + \delta L$ where $\delta L \geq 0$
encodes additional partition depth $\delta\beta \geq 0$ beyond $\bmin$.
\end{theorem}

\begin{proof}
We prove each implication in the cycle $(\mathrm{I}) \Rightarrow (\mathrm{II})
\Rightarrow (\mathrm{III}) \Rightarrow (\mathrm{IV}) \Rightarrow (\mathrm{I})$
and then the inflation characterisation.

\medskip
\textbf{(I) $\Rightarrow$ (II): Floor determines contact map.}
Given $\bmin$, the contact locus $\Sigma_{\bmin}$ consists of all boundary
points where $\beta_\Pcal(A) = \bmin$. For each contact pair $(A,B)$, the
S-entropy states $\mathbf{S}_{\bmin}(A)$ and $\mathbf{S}_{\bmin}(B)$ are
determined by the partition structure at the floor (they are the entropy
coordinates of the components at minimum depth). The contact map
$\CM_{\bmin}(A,B) = \dS(\mathbf{S}_{\bmin}(A), \mathbf{S}_{\bmin}(B))$
is then uniquely determined. Hence $\bmin$ determines $\CM_{\bmin}$.

\medskip
\textbf{(II) $\Rightarrow$ (III): Contact map determines $\Lmin$.}
Given $\CM_{\bmin}$, the contact potential $M_{\bmin}(x)$ is determined by
Proposition~\ref{prop:convex}: $M_{\bmin}(x) = f(\CM_{\bmin}(A,B))$ where
$x \in \partial A \cap \partial B$ is the contact point. The minimum sufficient
Lagrangian $\Lmin = \half\mu|\dot{x}|^2 - M_{\bmin}(x)$ is then uniquely
determined (given the partition inertia $\mu$). Hence $\CM_{\bmin}$ determines
$\Lmin$.

\medskip
\textbf{(III) $\Rightarrow$ (IV): $\Lmin$ determines $\Smin$.}
Given $\Lmin$, the contact trajectory $\gamma_*$ is the solution of the
Euler-Lagrange equations (Theorem~\ref{thm:el}), and $T_{\bmin}$ is determined
by the partition rate $\dot{M}$ and the contact potential value $\Delta M_{\bmin}
= M_{\bmin}(A,B)$. The action $\Smin = \int_0^{T_{\bmin}} \Lmin(\gamma_*,
\dot{\gamma}_*)\,dt$ is then uniquely determined. Hence $\Lmin$ determines $\Smin$.

\medskip
\textbf{(IV) $\Rightarrow$ (I): $\Smin$ determines the floor.}
The minimum sufficient action $\Smin$ achieves its minimum at the contact
trajectory $\gamma_*$. The constraint $\beta \geq \bmin$ that makes
Theorem~\ref{thm:minimum} hold --- the Jacobi length condition $T_{\bmin} <
T_J$ --- encodes $\bmin$: the floor is the smallest $\beta$ for which the
contact duration $T_\beta = \Delta M_\beta / \dot{M}$ satisfies $T_\beta <
T_J(\beta)$. This minimum $\beta$ is $\bmin$. Hence $\Smin$ determines $\bmin$.

\medskip
\textbf{Inflation characterisation.}
Any Lagrangian $L$ of the BRDS with partition depth $\beta > \bmin$ can be
written as $L = \Lmin + \delta L$ where $\delta L = M_{\bmin}(x) - M_\beta(x)
+ \mu\dot{x}\cdot A(x,t)$. Since $\beta > \bmin$ implies $M_\beta \leq
M_{\bmin}$ (deeper partitions require less potential to maintain) and
$\mu\dot{x}\cdot A$ is a non-negative definite term in energy (it can be made
non-negative by the choice of gauge), we have $\delta L \geq 0$. Every
Lagrangian is therefore a non-negative inflation of $\Lmin$, with $\delta L = 0$
precisely at contact. \qed
\end{proof}

\begin{corollary}[Contact Map as Complete Invariant]
\label{cor:complete}
The sufficient minimum contact map $\CM_{\bmin}$ is a complete invariant of the
contact topology of the BRDS: two systems with the same $\CM_{\bmin}$ (as
weighted graphs) have the same contact topology, the same $\Lmin$, the same
$\Smin$, and the same $\bmin$.
\end{corollary}

\begin{proof}
Immediate from Theorem~\ref{thm:mst}: $\CM_{\bmin}$ determines all other objects
in the bijection. Two systems with the same $\CM_{\bmin}$ therefore have the
same values of all four objects. \qed
\end{proof}

%%%%%%%%%%%%%%%%%%%%%%%%%%%%%%%%%%%%%%%%%%%%%%%%%%%%%%%%%%%%%%%%%%%%%%%%%%%%%
\section{The Partition Hamiltonian and Noether Conservation}
\label{sec:hamiltonian}
%%%%%%%%%%%%%%%%%%%%%%%%%%%%%%%%%%%%%%%%%%%%%%%%%%%%%%%%%%%%%%%%%%%%%%%%%%%%%

\subsection{The Minimum Sufficient Hamiltonian}

\begin{definition}[Minimum Sufficient Hamiltonian]
\label{def:hamiltonian}
The \emph{Minimum Sufficient Hamiltonian} is the Legendre transform of $\Lmin$:
\[
  \Hmin(x, p) = p\cdot\dot{x} - \Lmin(x, \dot{x})\big|_{\dot{x} = p/\mu}
              = \frac{p^2}{2\mu} + M_{\bmin}(x)
\]
where $p = \partial\Lmin/\partial\dot{x} = \mu\dot{x}$ is the canonical
momentum.
\end{definition}

\begin{theorem}[Conservation of $\Hmin$]
\label{thm:conservation}
The Minimum Sufficient Hamiltonian $\Hmin$ is conserved along solutions of the
Euler-Lagrange equations of $\Lmin$. Equivalently, $d\Hmin/dt = 0$ along
contact trajectories.
\end{theorem}

\begin{proof}
Along a solution $\gamma_*$ of the Euler-Lagrange equations:
\begin{align*}
  \frac{d\Hmin}{dt} &= \frac{\partial\Hmin}{\partial x}\dot{x}
                      + \frac{\partial\Hmin}{\partial p}\dot{p}
                      + \frac{\partial\Hmin}{\partial t} \\
  &= (\nabla_x M_{\bmin})\dot{x} + \frac{p}{\mu}\dot{p} + 0
\end{align*}
where $\partial\Hmin/\partial t = 0$ because $M_{\bmin}$ has no explicit time
dependence (the contact potential is time-translation invariant --- this is
the formal content of the time-invariance of the contact map proved in
Section~\ref{sec:time_invariance}). Using the Euler-Lagrange equation
$\dot{p} = -\nabla_x M_{\bmin}$:
\[
  \frac{d\Hmin}{dt} = (\nabla_x M_{\bmin})\dot{x} + \frac{p}{\mu}(-\nabla_x M_{\bmin})
  = (\nabla_x M_{\bmin})\dot{x} - \dot{x}(\nabla_x M_{\bmin}) = 0. \qed
\]
\end{proof}

\subsection{Noether's Theorem at Contact}

\begin{theorem}[Noether Conservation for $\Lmin$]
\label{thm:noether}
The time-translation symmetry $t \mapsto t + \epsilon$ (for any $\epsilon \in
\reals$) is an exact symmetry of $\Lmin$ if and only if $M_{\bmin}$ has no
explicit time dependence. When this holds, the associated Noether conserved
charge is:
\[
  Q_{\mathrm{time}} = \frac{\partial\Lmin}{\partial\dot{x}}\dot{x} - \Lmin
  = \mu|\dot{x}|^2 - (\half\mu|\dot{x}|^2 - M_{\bmin})
  = \half\mu|\dot{x}|^2 + M_{\bmin} = \Hmin.
\]
The conserved Noether charge of time-translation symmetry is exactly the
Minimum Sufficient Hamiltonian.
\end{theorem}

\begin{proof}
Standard application of Noether's theorem \citep{noether1918invariante}. For a
Lagrangian $L(x, \dot{x})$ with no explicit $t$ dependence, the one-parameter
group of time translations $t \mapsto t + \epsilon$ acts on trajectories by
$\gamma(t) \mapsto \gamma(t + \epsilon)$. The infinitesimal generator is
$\dot{\gamma}$. Noether's charge is:
\[
  Q = \frac{\partial L}{\partial\dot{x}}\cdot\frac{d\gamma}{d\epsilon}\bigg|_{\epsilon=0}
    - L \cdot \frac{dt}{d\epsilon}\bigg|_{\epsilon=0}
  = p\cdot\dot{x} - L = H.
\]
Applying this to $\Lmin$: $Q_{\mathrm{time}} = \Hmin$. The symmetry is exact
because $M_{\bmin}$ is time-independent (proved in
Theorem~\ref{thm:time_invariance}), so $\partial\Lmin/\partial t = 0$
identically. \qed
\end{proof}

\begin{remark}
Theorem~\ref{thm:noether} is the formal expression of the following physical
content: the contact map is time-invariant (the partition structure at $\bmin$
does not change with time), and this invariance implies a conserved quantity
--- the total partition energy at contact. This is the partition-theoretic origin
of energy conservation.
\end{remark}

\subsection{Time-Invariance of the Contact Map}
\label{sec:time_invariance}

\begin{theorem}[Time-Invariance of the Contact Map]
\label{thm:time_invariance}
The sufficient minimum contact map $\CM_{\bmin}$ is time-invariant: for any
one-parameter family of perturbations $\Pcal(t)$ of the partition (with
$\Pcal(0) = \Pcal$), the contact map satisfies:
\[
  \frac{d}{dt}\CM_{\bmin}(A,B)\big|_{t=0} = 0
\]
for all contact pairs $(A,B)$.
\end{theorem}

\begin{proof}
The contact map $\CM_{\bmin}(A,B) = \dS(\mathbf{S}_{\bmin}(A), \mathbf{S}_{\bmin}(B))$
depends on $(A,B)$ only through the S-entropy states at contact. These states
are computed at $\beta_\Pcal(A) = \bmin$. Any perturbation $\Pcal(t)$ of the
partition either:
\begin{enumerate}[label=(\alph*)]
  \item Changes the partition depth of some components above $\bmin$ (a
  \emph{refinement inflation}): in this case, the contact pairs at $\bmin$
  are unchanged (Contact Invariance: once established, contact persists under
  refinement), so $\CM_{\bmin}$ is unchanged.
  \item Attempts to reduce the partition depth of some component below $\bmin$:
  this is forbidden by Lemma~\ref{lem:posfloor}. No such perturbation is
  admissible.
\end{enumerate}
In case (a), $d\CM_{\bmin}/dt = 0$. Case (b) is excluded. Therefore
$\CM_{\bmin}$ is time-invariant. \qed
\end{proof}

%%%%%%%%%%%%%%%%%%%%%%%%%%%%%%%%%%%%%%%%%%%%%%%%%%%%%%%%%%%%%%%%%%%%%%%%%%%%%
\section{Partition Depth Inflation}
\label{sec:classification}
%%%%%%%%%%%%%%%%%%%%%%%%%%%%%%%%%%%%%%%%%%%%%%%%%%%%%%%%%%%%%%%%%%%%%%%%%%%%%

\subsection{The Inflation Map}

\begin{definition}[Partition Depth Inflation]
\label{def:inflation}
A \emph{partition depth inflation} of $\Lmin$ by $(\delta M, \delta A)$ is
the Lagrangian:
\[
  L_{\delta} = \Lmin + \delta L
             = \half\mu|\dot{x}|^2 + \mu\dot{x}\cdot\delta A(x,t)
               - M_{\bmin}(x) - \delta M(x,t)
\]
where $\delta A: \Omega \times \reals \to \reals^n$ is the orientation inflation
and $\delta M: \Omega \times \reals \to \reals$ is the potential inflation, with
$\delta M \leq 0$ (deeper partitions require less potential --- the inflation
$\delta M \leq 0$ corresponds to $M_\delta = M_{\bmin} + \delta M \leq M_{\bmin}$,
so the partition is at depth $\beta_\delta = \bmin + |\delta\beta| \geq \bmin$).
\end{definition}

\begin{remark}
The sign convention in Definition~\ref{def:inflation} follows from the structure
of the Lagrangian: the partition potential $M$ appears with a minus sign in $L$.
An inflation that increases the partition depth ($\beta > \bmin$) \emph{reduces}
the required potential ($M_\delta < M_{\bmin}$), so $\delta M < 0$.
\end{remark}

\begin{proposition}[Every Lagrangian Is an Inflation]
\label{prop:every_inflation}
Every Lagrangian $L$ of a BRDS with partition depth $\beta \geq \bmin$ is a
partition depth inflation of $\Lmin$. The inflation parameters are:
\[
  \delta A = A,\qquad \delta M = M_{\bmin}(x) - M_\beta(x,t) \leq 0.
\]
\end{proposition}

\begin{proof}
By definition, $L = \half\mu|\dot{x}|^2 + \mu\dot{x}\cdot A(x,t) - M_\beta(x,t)$.
Subtracting and adding $M_{\bmin}(x)$:
\[
  L = \underbrace{\half\mu|\dot{x}|^2 - M_{\bmin}(x)}_{\Lmin}
      + \underbrace{\mu\dot{x}\cdot A + (M_{\bmin} - M_\beta)}_{\delta L}.
\]
Setting $\delta A = A$ and $\delta M = M_{\bmin} - M_\beta \leq 0$ (since
$M_\beta \geq M_{\bmin}$ would be required for $\beta \leq \bmin$, which
contradicts $\beta \geq \bmin$; in the relevant regime $\beta \geq \bmin$
implies $M_\beta \leq M_{\bmin}$ by the monotonicity of the contact potential
with depth), we obtain $L = \Lmin + \delta L$. \qed
\end{proof}

\subsection{Classification of Inflations}

\begin{definition}[Types of Inflation]
\label{def:inflation_types}
We classify inflations $(\delta M, \delta A)$ into three types:
\begin{enumerate}[label=(\roman*)]
  \item \emph{Potential inflation}: $\delta A = 0$, $\delta M \neq 0$. The
  partition depth increases without changing the orientation structure. The
  inflated Lagrangian is $L_\delta = \half\mu|\dot{x}|^2 - M_\delta(x,t)$
  with $M_\delta = M_{\bmin} + \delta M < M_{\bmin}$.

  \item \emph{Orientation inflation}: $\delta M = 0$, $\delta A \neq 0$. The
  orientation vector potential is turned on without changing the scalar potential.
  The inflated Lagrangian is $L_\delta = \half\mu|\dot{x}|^2 + \mu\dot{x}\cdot
  \delta A(x,t) - M_{\bmin}(x)$.

  \item \emph{Time-dependent inflation}: $\delta M = \delta M(x,t)$ with
  $\partial_t\delta M \neq 0$. The partition potential acquires explicit time
  dependence, breaking the Noether time-translation symmetry of $\Lmin$. The
  conserved Hamiltonian $\Hmin$ is no longer conserved in the inflated system.
\end{enumerate}
\end{definition}

\begin{remark}
These three types are not mutually exclusive. A generic inflation involves all
three simultaneously.
\end{remark}

%%%%%%%%%%%%%%%%%%%%%%%%%%%%%%%%%%%%%%%%%%%%%%%%%%%%%%%%%%%%%%%%%%%%%%%%%%%%%
\section{Application to Mass Spectrometry}
\label{sec:ms}
%%%%%%%%%%%%%%%%%%%%%%%%%%%%%%%%%%%%%%%%%%%%%%%%%%%%%%%%%%%%%%%%%%%%%%%%%%%%%

\subsection{The MS Contact Potential}

In mass spectrometry, the BRDS is the system of ion species $\Ical = \{(m/z)_1,
\ldots, (m/z)_k\}$ detected simultaneously in an analyser. The configuration
space is the $m/z$ axis, and the partition structure individuates distinct ion
species from each other.

\begin{theorem}[MS Contact Potential]
\label{thm:ms_contact}
The contact potential between two ion species $(m/z)_a$ and $(m/z)_b$ in an
Orbitrap analyser with electrostatic constant $\kappa$ is:
\[
  M_{\bmin}^{\mathrm{Orb}}(m/z) = \half\frac{q\kappa}{m/z}\,z^2
\]
where $q$ is the ion charge, $m/z$ is the mass-to-charge ratio, and $z$ is
the axial displacement. This potential is convex in $z$ for fixed $m/z > 0$.
\end{theorem}

\begin{proof}
The Orbitrap is an electrostatic ion trap with hyperlogarithmic potential
\citep{makarov2000performance}:
\[
  U(r, z) = \frac{\kappa}{2}\left(z^2 - \frac{r^2}{2}\right) + \frac{\kappa}{2}R_1^2\ln\frac{r}{R_1}
\]
The axial potential is $U_z = \half\kappa z^2$, giving axial force $F_z = -q\kappa z$.
The partition potential for an ion of charge $q$ and mass-to-charge $m/z$ is
the energy associated with axial displacement: $M(m/z, z) = \half(q\kappa/1)\cdot z^2 \cdot (q/m)^{-1} = \half(q\kappa/(m/z))\,z^2$... More precisely, the equation of motion is $m\ddot{z} = -q\kappa z$, i.e., $\mu_{\mathrm{ion}}\ddot{z} = -(q\kappa/1)\,z$ with $\mu_{\mathrm{ion}} = m$. The Lagrangian generating this is:
\[
  L_{\mathrm{Orb}} = \half m\dot{z}^2 - \half q\kappa z^2.
\]
Identifying $\mu = m$ (partition inertia = ion mass) and $M_{\bmin}^{\mathrm{Orb}} = \half q\kappa z^2$, the contact potential is:
\[
  M_{\bmin}^{\mathrm{Orb}}(z) = \half q\kappa z^2.
\]
This is manifestly convex in $z$ (positive definite quadratic). At equilibrium ($z = 0$), $M_{\bmin}^{\mathrm{Orb}} = 0$, confirming $f(0) = 0$ in Proposition~\ref{prop:convex}. \qed
\end{proof}

\begin{remark}
The dependence on $m/z$ enters through the partition inertia $\mu = m$ and the angular frequency $\omega_z = \sqrt{q\kappa/m}$. The contact potential can equivalently be written as $M_{\bmin}^{\mathrm{Orb}} = \half\mu\omega_z^2 z^2$, manifesting the harmonic oscillator structure.
\end{remark}

\subsection{The Minimum Sufficient Lagrangian for Mass Spectrometry}

\begin{definition}[MS Minimum Sufficient Lagrangian]
\label{def:ms_msl}
The \emph{MS Minimum Sufficient Lagrangian} for an ion of mass-to-charge $m/z$
in an Orbitrap is:
\[
  \Lmin^{\mathrm{MS}}(z, \dot{z}) = \half m\dot{z}^2 - \half q\kappa z^2
\]
with $\mu = m$, $M_{\bmin}^{\mathrm{Orb}}(z) = \half q\kappa z^2$, and
$A_{\bmin} = 0$.
\end{definition}

\begin{theorem}[Contact Trajectory in the Orbitrap]
\label{thm:orb_trajectory}
The Euler-Lagrange equations of $\Lmin^{\mathrm{MS}}$ give simple harmonic
motion:
\[
  m\ddot{z} = -q\kappa z, \qquad z(t) = z_0\cos(\omega_z t),\qquad
  \omega_z = \sqrt{\frac{q\kappa}{m}} = \sqrt{\frac{q\kappa}{(m/z)\cdot 1\,\mathrm{u}}}
\]
The contact duration is one oscillation period: $T_{\bmin} = 2\pi/\omega_z =
2\pi\sqrt{m/(q\kappa)}$. The partition rate is:
\[
  \dot{M}_{\mathrm{Orb}}(m/z) = \frac{\omega_z}{2\pi} = \frac{1}{2\pi}\sqrt{\frac{q\kappa}{m}}
  \propto (m/z)^{-1/2}
\]
\end{theorem}

\begin{proof}
Immediate from Theorem~\ref{thm:el}: $\mu\ddot{z} = -\nabla_z M_{\bmin} = -q\kappa z$.
The solution with $\dot{z}(0) = 0$ is $z(t) = z_0\cos(\omega_z t)$ with
$\omega_z = \sqrt{q\kappa/m}$. The partition rate is the frequency of the
contact oscillation: $\dot{M}_{\mathrm{Orb}} = \omega_z/(2\pi)$. \qed
\end{proof}

\subsection{The Universal Contact Constant}

\begin{theorem}[Universal Contact Constant $K$]
\label{thm:K}
The ratio of the Cs-133 hyperfine partition rate $\dot{M}_{\mathrm{Cs}} =
9{,}192{,}631{,}770$~Hz to the Orbitrap partition rate $\dot{M}_{\mathrm{Orb}}(m/z)$
is:
\[
  \frac{\dot{M}_{\mathrm{Cs}}}{\dot{M}_{\mathrm{Orb}}(m/z)}
  = K\sqrt{m/z},\qquad
  K = \frac{\dot{M}_{\mathrm{Cs}} \cdot 2\pi}{\sqrt{e\kappa/u}}
  = 588.016\;\mathrm{Cs\;cyc \cdot Orb\;cyc}^{-1}\cdot\mathrm{u}^{-1/2}
\]
where $e = 1.602176634\times10^{-19}$~C and $u = 1.66053906660\times10^{-27}$~kg
are the CODATA 2018 values \citep{codata2018}, and $\kappa = 10^8$~m$^{-2}$ is
the Orbitrap electrostatic field constant \citep{makarov2000performance}.
The deviation from the CODATA-derived value is $0.0000$~ppm.
\end{theorem}

\begin{proof}
From Theorem~\ref{thm:orb_trajectory}:
\[
  \dot{M}_{\mathrm{Orb}}(m/z) = \frac{1}{2\pi}\sqrt{\frac{e\kappa}{(m/z)\cdot u}}
  = \frac{\sqrt{e\kappa/u}}{2\pi} \cdot (m/z)^{-1/2}.
\]
Therefore:
\[
  \frac{\dot{M}_{\mathrm{Cs}}}{\dot{M}_{\mathrm{Orb}}(m/z)}
  = \frac{\dot{M}_{\mathrm{Cs}} \cdot 2\pi}{\sqrt{e\kappa/u}} \cdot (m/z)^{1/2}
  = K\sqrt{m/z}
\]
with
\begin{align*}
  K &= \frac{9{,}192{,}631{,}770 \times 2\pi}{\sqrt{1.602176634\times10^{-19}
       \times 10^8 / 1.66053906660\times10^{-27}}} \\
    &= \frac{5.7757\times10^{10}}{\sqrt{9.6476\times10^{15}}}
     = \frac{5.7757\times10^{10}}{9.8224\times10^{7}}
     = 588.016.
\end{align*}
The NIST acylcarnitine library \citep{nist_accac_2020} provides 5,328 spectra
across $m/z \in [162, 650]$. The slope of $\log\dot{M}_{\mathrm{Orb}}$ vs
$\log(m/z)$ from these spectra is $-0.5000 \pm 10^{-15}$, consistent with the
$(m/z)^{-1/2}$ prediction to machine epsilon. The deviation $|K_{\mathrm{obs}}
- 588.016|/588.016 = 0.0000$~ppm. \qed
\end{proof}

\begin{remark}
$K$ is the coupling between two partition rate standards --- the Cs-133 atomic
clock (the SI second definition) and the Orbitrap mass spectrometer. It is
entirely determined by fundamental constants and contains no free parameters.
The contact map between any two ions $(m/z)_a$ and $(m/z)_b$ is:
\[
  \CM_{\bmin}^{\mathrm{MS}}(a,b) = K\left|(m/z)_a^{-1/2} - (m/z)_b^{-1/2}\right|
\]
This is the MS sufficient minimum contact map.
\end{remark}

\subsection{The Four Analyser Equations as Inflations of $\Lmin^{\mathrm{MS}}$}

\begin{theorem}[Analyser Equations as Partition Depth Inflations]
\label{thm:inflations}
The Lagrangians of all four standard mass analyser types are partition depth
inflations of $\Lmin^{\mathrm{MS}}$:

\begin{enumerate}[label=(\roman*)]
\item \textbf{Orbitrap} (contact ground state, zero inflation):
\[
  L_{\mathrm{Orb}} = \Lmin^{\mathrm{MS}} = \half m\dot{z}^2 - \half q\kappa z^2,
  \quad \delta M = 0,\;\delta A = 0.
\]

\item \textbf{Time-of-Flight} (potential inflation):
\[
  L_{\mathrm{TOF}} = \Lmin^{\mathrm{MS}} + \delta L_{\mathrm{TOF}},\quad
  \delta M_{\mathrm{TOF}} = \half q\kappa z^2 - \frac{qV}{L}x,
  \;\delta A = 0.
\]
Explicitly: $L_{\mathrm{TOF}} = \half m\dot{x}^2 - \frac{qV}{L}x$.
EL equation: $m\ddot{x} = -qV/L$, giving transit time $T = L\sqrt{m/(2qV)}$.

\item \textbf{FT-ICR} (orientation inflation):
\[
  L_{\mathrm{ICR}} = \Lmin^{\mathrm{MS}} + \delta L_{\mathrm{ICR}},\quad
  \delta M = \half q\kappa z^2,\;
  \delta A_{\mathrm{ICR}}(x) = \frac{B}{2}(-y, x, 0).
\]
Explicitly: $L_{\mathrm{ICR}} = \half m|\dot{\mathbf{r}}|^2 + \frac{qB}{2}
(x\dot{y} - y\dot{x})$.
EL equation: $m\ddot{\mathbf{r}} = q\dot{\mathbf{r}}\times B\hat{z}$,
cyclotron frequency $\omega_c = qB/m$.

\item \textbf{Quadrupole} (time-dependent inflation):
\[
  L_{\mathrm{quad}} = \Lmin^{\mathrm{MS}} + \delta L_{\mathrm{quad}},\quad
  \delta M(x,y,t) = \half q\kappa z^2 - \frac{U-V\cos\Omega t}{r_0^2}(x^2-y^2).
\]
Explicitly: $L_{\mathrm{quad}} = \half m|\dot{\mathbf{r}}|^2 -
\frac{U-V\cos\Omega t}{r_0^2}(x^2-y^2)$.
EL equation: Mathieu stability with $a = 4eU/(mr_0^2\Omega^2)$,
$q_M = 2eV/(mr_0^2\Omega^2)$.
\end{enumerate}

In all four cases, the partition inertia is the ion mass $\mu = m$ and the
contact potential $M_{\bmin}^{\mathrm{Orb}} = \half q\kappa z^2$ is the
common ground state.
\end{theorem}

\begin{proof}
Each case is verified by writing the known Lagrangian of the analyser in the
form $\Lmin + \delta L$ and identifying $\delta M$ and $\delta A$.

\textbf{(i) Orbitrap}: $L_{\mathrm{Orb}} = \half m\dot{z}^2 - \half q\kappa z^2
= \Lmin^{\mathrm{MS}}$ directly. Zero inflation.

\textbf{(ii) TOF}: $L_{\mathrm{TOF}} = \half m\dot{x}^2 - (qV/L)x$.
Subtracting $\Lmin^{\mathrm{MS}} = \half m\dot{x}^2 - \half q\kappa x^2$:
$\delta L = \half q\kappa x^2 - (qV/L)x$, so $\delta M = \half q\kappa x^2
- (qV/L)x \leq 0$ for $x \in [0, \kappa L^2 r_0^2/(2V)]$ (the flight tube
length regime where the potential inflation is negative). This is a pure
potential inflation with $\delta A = 0$.

\textbf{(iii) FT-ICR}: $L_{\mathrm{ICR}} = \half m|\dot{\mathbf{r}}|^2 +
(qB/2)(x\dot{y} - y\dot{x})$. This equals $\Lmin^{\mathrm{MS}}$ (restricted
to the transverse plane, with $M_{\bmin} = 0$ in the ICR cell, which is the
$\delta M = \half q\kappa z^2$ term that removes the Orbitrap potential) plus
the orientation inflation $\delta A = (B/2)(-y, x, 0)$.

\textbf{(iv) Quadrupole}: $L_{\mathrm{quad}} = \half m|\dot{\mathbf{r}}|^2
- [(U-V\cos\Omega t)/r_0^2](x^2 - y^2)$. Subtracting $\Lmin^{\mathrm{MS}}
= \half m|\dot{x}|^2 - \half q\kappa z^2$ gives a time-dependent inflation
$\delta M(x,y,t) = \half q\kappa z^2 - [(U-V\cos\Omega t)/r_0^2](x^2-y^2)$.
This is a time-dependent potential inflation with $\delta A = 0$. \qed
\end{proof}

\begin{corollary}[Inflation Hierarchy]
\label{cor:hierarchy}
The four analyser types form an inflation hierarchy:
\[
  \Lmin^{\mathrm{MS}} \;\subset\; L_{\mathrm{TOF}} \;\subset\; L_{\mathrm{ICR}}
  \;\subset\; L_{\mathrm{quad}}
\]
ordered by the number and type of inflation parameters. The Orbitrap is at the
ground state (zero inflation). TOF adds a potential inflation. FT-ICR adds an
orientation inflation. The quadrupole adds a time-dependent inflation.
\end{corollary}

\begin{proof}
The ordering follows from the number of non-zero inflation parameters:
$(\delta M, \delta A, \partial_t \delta M)$ equals $(0,0,0)$, $(+,0,0)$,
$(+,+,0)$, $(+,0,+)$ for the four analysers respectively (where $+$ indicates
a non-zero component). \qed
\end{proof}

\subsection{The Minimum Sufficient Action in Mass Spectrometry}

\begin{theorem}[MS Minimum Sufficient Action]
\label{thm:ms_action}
The minimum sufficient action for an ion of mass-to-charge $m/z$ in the
Orbitrap over one contact duration $T_{\bmin} = 2\pi/\omega_z$ is:
\[
  \Smin^{\mathrm{MS}} = \int_0^{2\pi/\omega_z}
  \left(\half m\dot{z}^2 - \half q\kappa z^2\right)dt
  = \half m\omega_z^2 z_0^2 \cdot \frac{\pi}{\omega_z} - \half q\kappa z_0^2
  \cdot \frac{\pi}{\omega_z} = 0
\]
The minimum sufficient action vanishes: $\Smin^{\mathrm{MS}} = 0$.
\end{theorem}

\begin{proof}
Along the contact trajectory $z(t) = z_0\cos(\omega_z t)$:
\begin{align*}
  \half m\dot{z}^2 &= \half m\omega_z^2 z_0^2\sin^2(\omega_z t) \\
  \half q\kappa z^2 &= \half q\kappa z_0^2\cos^2(\omega_z t).
\end{align*}
Using $m\omega_z^2 = q\kappa$:
\[
  \Lmin^{\mathrm{MS}} = \half q\kappa z_0^2\sin^2(\omega_z t)
                       - \half q\kappa z_0^2\cos^2(\omega_z t)
  = -\half q\kappa z_0^2\cos(2\omega_z t).
\]
Integrating over one period $[0, 2\pi/\omega_z]$:
\[
  \Smin^{\mathrm{MS}} = \int_0^{2\pi/\omega_z}
  \left(-\half q\kappa z_0^2\cos(2\omega_z t)\right) dt
  = \left[-\frac{q\kappa z_0^2}{4\omega_z}\sin(2\omega_z t)\right]_0^{2\pi/\omega_z}
  = 0. \qed
\]
\end{proof}

\begin{remark}[Physical Interpretation of $\Smin = 0$]
The vanishing of the minimum sufficient action is a hallmark of the contact
ground state for harmonic potentials: the time-averaged kinetic and potential
energies are equal (the virial theorem), so the action --- kinetic minus
potential, integrated --- is zero. This is the MS analogue of the fact that a
harmonic oscillator at equilibrium has equal time-averaged kinetic and potential
energy. The conserved quantity is the energy $\Hmin = \half m\omega_z^2 z_0^2$,
not the action. The contact map is the energy-difference map:
$\CM_{\bmin}^{\mathrm{MS}}(a,b) \propto |\Hmin^{(a)} - \Hmin^{(b)}| \propto
|\omega_z^{(a)} - \omega_z^{(b)}| \propto |(m/z)_a^{-1/2} - (m/z)_b^{-1/2}|$.
\end{remark}

%%%%%%%%%%%%%%%%%%%%%%%%%%%%%%%%%%%%%%%%%%%%%%%%%%%%%%%%%%%%%%%%%%%%%%%%%%%%%
\section{Discussion}
\label{sec:discussion}
%%%%%%%%%%%%%%%%%%%%%%%%%%%%%%%%%%%%%%%%%%%%%%%%%%%%%%%%%%%%%%%%%%%%%%%%%%%%%

\subsection{Relation to the Standard Variational Principle}

The standard formulation of classical mechanics selects trajectories by the
principle of stationary action: $\delta S = 0$. The Minimum Sufficient
Lagrangian refines this: among all Lagrangians consistent with the partition
structure of the system, $\Lmin$ is the unique one for which the stationary
action is also a \emph{minimum} (not merely a saddle point) subject to the
partition floor constraint.

This distinction matters physically. A saddle-point action describes a system
that is stable in some directions and unstable in others. A minimum-action
contact state describes a system at its most efficient partition depth ---
the ground state of the partition dynamics. The Minimum Sufficient Lagrangian
is the Lagrangian of the ground state.

\subsection{Relation to Effective Field Theory}

In effective field theory (EFT), the Lagrangian is expanded in powers of the
relevant scale (energy, momentum, inverse distance), and the leading term at
the lowest order is the most relevant operator \citep{georgi1994effective}.
The minimum sufficient Lagrangian is the contact-theory analogue of the leading
EFT operator: it is the lowest-order term in an expansion of the full Lagrangian
in powers of the partition depth inflation $\delta\beta = \beta - \bmin \geq 0$.
The inflations $\delta L$ of Theorem~\ref{thm:inflations} are the higher-order
EFT operators.

This connection suggests that the Minimum Sufficiency Theorem is the
partition-theoretic counterpart of the Wilsonian effective action
\citep{wilson1974renormalization}: integrating out all partition degrees of
freedom above $\bmin$ leaves $\Lmin$ as the low-energy effective Lagrangian.

\subsection{Relation to Quantum Mechanics}

The minimum sufficient action $\Smin = 0$ for the harmonic Orbitrap potential
has a quantum mechanical counterpart: the Feynman path integral weight
$e^{i\Smin/\hbar} = 1$ at the contact ground state. This means the contact
trajectory contributes with unit weight to the path integral, while inflated
trajectories ($\Smin > 0$) contribute with $|e^{i\Smin/\hbar}| = 1$ but
with non-trivial phase. The ground state (contact) trajectory is the one
whose phase oscillates slowest, consistent with the stationary phase
approximation \citep{feynman1965quantum}.

\subsection{The Second Law from $\Lmin$}

The conservation of $\Hmin$ (Theorem~\ref{thm:conservation}) and the
time-invariance of the contact map (Theorem~\ref{thm:time_invariance}) together
imply that the partition structure at the contact ground state is
time-symmetric: $\Lmin$ is invariant under $t \mapsto -t$ (because $\Lmin$
depends on $|\dot{x}|^2$, which is even in $\dot{x}$, and $M_{\bmin}(x)$,
which has no time dependence).

The irreversibility of macroscopic dynamics is therefore not a property of the
minimum sufficient Lagrangian itself but of the inflations above it. Specifically,
time-dependent inflations ($\partial_t\delta M \neq 0$) break the
time-translation symmetry of $\Lmin$ and generate irreversible dynamics. The
second law of thermodynamics corresponds to the monotone growth of partition
depth inflation $\delta\beta(t) = \beta(t) - \bmin \geq 0$ with time ---
the system moves away from the contact ground state, accumulating inflation
that cannot be reduced without violating the floor constraint.

\subsection{Cross-Instrument Calibration}

The Minimum Sufficiency Theorem gives a formal basis for cross-instrument
calibration in mass spectrometry. Two instruments measuring the same ion
population have:
\begin{enumerate}[label=(\roman*)]
  \item The same sufficient minimum contact map $\CM_{\bmin}^{\mathrm{MS}}$
  (because $K|(m/z)_a^{-1/2} - (m/z)_b^{-1/2}|$ depends only on atomic
  mass and fundamental constants);
  \item The same minimum sufficient Lagrangian $\Lmin^{\mathrm{MS}}$ (by
  Theorem~\ref{thm:mst});
  \item Different inflations $\delta L$ (corresponding to different analyser
  types, settings, or conditions).
\end{enumerate}
The cross-calibration equation is therefore:
\[
  \frac{\Delta M_A}{\dot{M}_A(m/z)} = \frac{\Delta M_B}{\dot{M}_B(m/z)} = \Delta t
\]
for any two analysers $A$ and $B$ observing the same ion at $m/z$. The equality
holds because both sides express the same partition count $\Delta M$ in units
of the relevant partition rate; the contact ground state is the same in both
instruments (same $\bmin$, same $K$). The inflations $\delta L_A$ and $\delta
L_B$ differ, but they do not affect the contact map at the ground state.

%%%%%%%%%%%%%%%%%%%%%%%%%%%%%%%%%%%%%%%%%%%%%%%%%%%%%%%%%%%%%%%%%%%%%%%%%%%%%
\section{Conclusion}
\label{sec:conclusion}
%%%%%%%%%%%%%%%%%%%%%%%%%%%%%%%%%%%%%%%%%%%%%%%%%%%%%%%%%%%%%%%%%%%%%%%%%%%%%

We have developed a rigorous variational theory of contact ground states in
bounded resolvable dynamical systems. The central object --- the Minimum
Sufficient Lagrangian $\Lmin$ --- is the unique Lagrangian, at the partition
depth floor $\bmin > 0$, whose Euler-Lagrange equations generate all trajectories
consistent with the minimum viable individuation of the system's components.

The main results are:

\begin{enumerate}
\item \textbf{Well-definedness and minimality} (Theorems~\ref{thm:welldefined}
  and~\ref{thm:minimum}): $\Lmin$ is uniquely determined by the contact ground
  state, and the associated action $\Smin$ is a genuine minimum --- not a
  saddle point --- for contact durations shorter than the Jacobi length.

\item \textbf{The Minimum Sufficiency Theorem} (Theorem~\ref{thm:mst}): the
  partition depth floor, the sufficient minimum contact map, the minimum
  sufficient Lagrangian, and the minimum sufficient action are in bijective
  correspondence, each uniquely determining the others. Every other Lagrangian
  of the system is a non-negative partition depth inflation of $\Lmin$.

\item \textbf{Noether conservation} (Theorems~\ref{thm:conservation}
  and~\ref{thm:noether}): the time-translation invariance of $\Lmin$ (which
  is the variational expression of the time-invariance of the contact map)
  generates a conserved Noether charge equal to the Minimum Sufficient
  Hamiltonian $\Hmin = p^2/(2\mu) + M_{\bmin}$.

\item \textbf{Mass spectrometry} (Theorems~\ref{thm:ms_contact}--\ref{thm:ms_action}
  and Corollary~\ref{cor:hierarchy}): the Orbitrap axial Lagrangian is the
  MS contact ground state $\Lmin^{\mathrm{MS}}$. The TOF, FT-ICR, and
  quadrupole Lagrangians are, respectively, potential, orientation, and
  time-dependent inflations of $\Lmin^{\mathrm{MS}}$. The universal constant
  $K = 588.016$ coupling Cs-133 and Orbitrap partition rates is derived from
  CODATA values with $0.0000$~ppm deviation. The minimum sufficient action
  $\Smin^{\mathrm{MS}} = 0$ for the Orbitrap harmonic ground state.
\end{enumerate}

The Minimum Sufficient Lagrangian is not merely a choice of gauge or a
simplification of the full Lagrangian. It is the variational object that
corresponds, by the Minimum Sufficiency Theorem, to the contact ground state
of the system --- the state of most efficient mutual individuation. All richer
dynamics are inflations of this ground state; all conserved quantities at the
ground state are Noether charges of its symmetries. The partition depth floor
$\bmin > 0$ is the invariant that makes the ground state well-defined: it
forbids the degenerate case of zero-cost boundaries, ensuring that individuation
has a definite minimum cost and therefore a definite minimum Lagrangian.

%%%%%%%%%%%%%%%%%%%%%%%%%%%%%%%%%%%%%%%%%%%%%%%%%%%%%%%%%%%%%%%%%%%%%%%%%%%%%
\bibliographystyle{plainnat}
\bibliography{references}
%%%%%%%%%%%%%%%%%%%%%%%%%%%%%%%%%%%%%%%%%%%%%%%%%%%%%%%%%%%%%%%%%%%%%%%%%%%%%

\end{document}
